\documentclass{beamer}

\beamertemplatenavigationsymbolsempty

\usepackage{etex}
% packages required  by CVMono
\usepackage{mathptmx}
\usepackage{helvet}
\usepackage{courier}
\usepackage{type1cm}
\usepackage{makeidx}
\usepackage{graphicx}
\usepackage{multicol}
\usepackage[bottom]{footmisc}
\usepackage{placeins}



%%%
\usepackage[UTF8x]{inputenc}
\usepackage[german,english]{babel}
\usepackage{amsmath}
\usepackage{amssymb}
%\usepackage{amsthm}
\usepackage{pifont}% for checkmarks \ding{51}
\makeatletter
\@ifclassloaded{beamer}{}{\usepackage{paralist}}
\makeatother
\usepackage{float}
\usepackage{adjustbox}

\usepackage[%
    font={small,sf},
    labelfont=bf,
    format=hang,    
    format=plain,
    margin=0pt,
    width=0.8\textwidth,
]{caption}
\usepackage[list=true]{subcaption}
%\usepackage{alltt}

% symbols for the naturals (\mathbbm{N}), integers (\mathbbm{Z}), etc.
\usepackage{bbm}

% nice tables (read the accompanying docs!)
\usepackage{booktabs}
\usepackage{hyperref}% for bookmarks in PDF file
\usepackage{url}%showing url in bibtex
% source code and diagrams
\usepackage{tikz}
\usepackage{tikz-timing}

% code and monostace
\usepackage{listings}
\usepackage{fancyvrb}
\usepackage{fixltx2e}
\usepackage{textcomp}

\usepackage{datetime}



\usepackage{tikz}
\usepackage{import}

% extract lemmas, inserted lemmas
\usepackage{verbatimcopy}
%% "touch foo.ext" before first use
\VerbatimCopy{lem_def_summary}{lem_def_summary_final}

% inference rules
\usepackage{mathpartir}

\let\proof\relax
\let\endproof\relax
\usepackage{amsthm}

% set literals, \Set{ ... | ... }, \Set{ ..., ..., ..., ... }
\usepackage{braket}

% interval literals, \interval{...}{...}, \interval[open left]{...}{...} etc
\usepackage{interval}
\intervalconfig{
	soft  open  fences,
	separator symbol=:,
}

\usepackage{import}
\usepackage{bm}

% \centernot : like \not, but better alignment
\usepackage{centernot}

% has nice things like \cmark and \xmark
\usepackage{pifont}
\usepackage{wasysym}


\usepackage{multirow}
\usepackage{nameref}

% has multiple references grouped together in clever ways; \cref
\usepackage[capitalise]{cleveref}

% manipulates strings
\usepackage{xstring}
%\theoremstyle{definition}
%\newtheorem{definition}{Definition}
%\newtheorem{lemma}{Lemma}
%\newtheorem{theorem}{Theorem}

% note: the next 2 environments are used only in sysbook; have to be adapted wrt to the style reference
\newtheorem{invariant}{Invariant}
\newtheorem{ccond}{Context Condition}

% make float placement on top the default:
\floatplacement{figure}{t}
\floatplacement{table}{t}
\renewcommand{\floatpagefraction}{0.8}

% tikz
\tikzset{timing/slope=0.1}
\tikzset{timing/.style = {y=2.5ex, x = 4.0ex}}

% other
\graphicspath{{figures/}}
\bibliographystyle{alpha}

% remove later
\newcommand{\typo}[1]{\textcolor{blue}{\textbf{#1}}}
\newcommand{\bbox}[1]{\fcolorbox{blue}{white}{#1}}

%%%%%%%%%%%%%%%%%%%%%%%%%%%%%%%%%%%%%%%%%%%%%%%%%%%%%%%%%%%%%%%%
% Define your desired shortcuts here:
\mathchardef\hy="2D
\newcommand{\cm}{\mbox{,}}
\newcommand{\setB}{\mathbb{B}}
\newcommand{\setN}{\mathbb{N}}
\newcommand{\setZ}{\mathbb{Z}}
\newcommand{\setR}{\mathbb{R}}
\newcommand{\setS}{\mathbb{S}}
\newcommand{\B}{\mathbb{B}}
\newcommand{\N}{\mathbb{N}}
\newcommand{\R}{\mathbb{R}}
\newcommand{\nat}{\mathbb{N}}
\newcommand{\Z}{\mathbb{Z}}
\newcommand{\NOT}[1]{\overline{#1}}
\newcommand{\AND}{\land}
\newcommand{\OR}{\lor}
\newcommand{\Sum}{\sum\limits}
\newcommand{\tmod}{\mathrel{\mathrm{tmod}}}
\newcommand{\pskip}{\smallskip\par\noindent}
\newcommand{\ignore}[1]{\relax}
\newcommand{\figscale}{0.834}
%\renewcommand{\pskip}{\smallskip\bbox{\mbox{\quad}}\par\noindent}
\newcommand{\smalltt}[1]{\text{\small\texttt{#1}}}
\newcommand{\super}[1]{\textsuperscript{#1}}
\newcommand{\sub}[1]{\textsubscript{#1}}
\newcommand{\langlett}{{\fontfamily{cmbr}\selectfont\textlangle}}
\newcommand{\ranglett}{{\fontfamily{cmbr}\selectfont\textrangle}}
\renewcommand{\theFancyVerbLine}{\text{\small\arabic{FancyVerbLine}:}}
\renewenvironment{verbatim}[0]{\Verbatim}{\endVerbatim}

% figures
\newcommand{\tfig}[3]{%
\centering%
\adjustbox{scale = \figscale}{\input{figures/#1.pdf_t}}%
\caption{#3}\label{fig:#2}%
}

\newcommand{\tsubfigV}[3]{%
\subcaptionbox{\label{#2} #3}{\import{figures/}{#1.pdf_tex}}%
}

\newcommand{\includefig}[1]{\import{figures/}{#1.pdf_tex}}


\newcommand{\tschemfig}[4]{%
\figure[tp]%
\addtolength{\subfigcapskip}{0.1in}%
\centering%
\subfigure[\label{fig:#1} symbol]{\trimbox{-0.5\textwidth+0.5\width 0pt 0pt}{\raisebox{#3mm}{\adjustbox{scale = \figscale}{\input{figures/#1.pdf_t}}}}}\\%
\subfigure[\label{fig:#1impl} implementation]{\trimbox{-0.5\textwidth+0.5\width 0pt 0pt}{\raisebox{#4mm}{\adjustbox{scale = \figscale}{\input{figures/#1impl.pdf_t}}}}}%
\caption{#2}\label{fig:#1-all}%
\endfigure%
}

%note: command "remark" is already defined in the style file
%\newcommand{\remark}[1]{\marginpar{\framebox{\parbox[t]{16mm}{
  %\tiny\raggedright #1}}}}
%\newcommand{\oldremark}[1]{\relax}

% code
\fvset{
	frame=lines,
	framerule=0.2pt,
	framesep=4pt,
	resetmargins=true,
	xleftmargin=16pt,
	xrightmargin=16pt,
%	numbers=left,
%	numberblanklines=false,
%	numbersep=-16pt,
	tabsize=0,
	fontfamily=courier,
	fontsize=\small,
	commandchars=\\\@\#
}

\lstnewenvironment{code}
    {\lstset{}%
      \csname lst@SetFirstLabel\endcsname}
    {\csname lst@SaveFirstLabel\endcsname}
    \lstset{
      language=Haskell,
      mathescape = true,
      commentstyle= \sffamily\itshape,
      basicstyle=\small\ttfamily,
      keywordstyle=\color{DarkBlue}\bfseries,
      flexiblecolumns=false,
      basewidth={0.5em,0.45em},
      morecomment=[l]{//},
      numbers = left,
      numberstyle = \color{Gray}\tiny\ttfamily,
      breaklines = true,
      literate={+}{{$+$}}1 {*}{{$*$}}1 {=}{{$=$}}1 {!=}{{$\neq$}}2
               {>}{{$>$}}1 {<}{{$<$}}1 {\\}{{$\lambda$}}1
               {->}{{$\rightarrow$}}2 {>=}{{$\geq$}}2 {<-}{{$\leftarrow$}}2
               {<=}{{$\leq$}}2 {=>}{{$\Rightarrow$}}2
               {>>}{{$\gg$}}1
               {<<}{{$\ll$}}1
               {|}{{$\mid$}}1
               {||}{{$\vee$}}2 {&&}{{$\wedge$}}2
               {@}{{$\circ$}}1
               {?}{{{\color{DarkBlue}{?}}}}1
               {...}{{$\ldots$}}3
    }

% forbids footnotes to travel
%\interfootnotelinepenalty=1000000000
\interfootnotelinepenalty=10000

% theorems of amsthm

% config of the tikz package
\usetikzlibrary{automata,positioning}

%new shorthands

% Tombstone at end of proof (amsthm does this now)
%\renewcommand\endproof{~\hfill\qed}

% std mathcal font
\DeclareMathAlphabet{\mathcal}{OMS}{cmsy}{m}{n}

% new utils
\newcommand{\tikscale}{1.0}

\newcommand{\tbl}[3]{%
\caption{#3}%
\label{tab:#2}%
\centering%
\input{tables/#1}%
}

\newcommand{\tpic}[3]{%
\centering%
\adjustbox{scale = \tikscale}{%
\tikzpicture[%
	>=stealth,%
	shorten >=0.5pt,%
	node distance=2.5cm,%
	on grid,%
	auto,%
	every node/.append style={very thin},%
	every path/.append style={very thin}%
]%
\input{drawings/#1}%
\endtikzpicture}%
\caption{#3}\label{fig:#2}%
}

\newcommand{\tfigscaled}[4]{%
\centering%
\adjustbox{scale = #4}{\input{figures/#1.pdf_t}}%
\caption{#3}\label{fig:#2}%
}

%%% REDUCTION PROOFS
%% units making steps
\newcommand{\MMU}{\mathrm{MMU}}
\newcommand{\SB}{\mathrm{SB}}
\newcommand{\APIC}{\mathrm{APIC}}
\newcommand{\IOAPIC}{\mathrm{IOAPIC}}
\newcommand{\DEV}{\mathrm{DEV}}

%% classes of schedules
\newcommand{\LSBS}{\mathrm{LSBS}}
\newcommand{\ABS}{\mathrm{ABS}}
\newcommand{\REDUCIBLE}{\mathrm{RED}}
\newcommand{\IND}{\mathrm{ORD}}

\newcommand{\IBLOCK}{\mathrm{IBLOCK}}
\newcommand{\CBLOCK}{\mathrm{CBLOCK}}
\newcommand{\EIBLOCK}{\mathrm{EIBLOCK}}
%% 
\newcommand{\sbcomp}[1]{{\overline{#1}}}


%% Makes it into a single word...
\newcommand{\inputsafe}{in\hy safe}
%% communication relations
\DeclareMathOperator{\fwdsync}\vartriangleright
\DeclareMathOperator{\rsync}{{\ooalign{$\vartriangleright$\cr\kern0.04em
  $\ast$\cr}}}
\DeclareMathOperator{\sync}{\blacktriangleright}
\DeclareMathOperator{\talks}{\sync\!\!\!\!\!\fwdsync}
\DeclareMathOperator{\rtalks}{\sync\!\!\!\!\!\rsync}
\DeclareMathOperator{\backsync}{\vartriangleleft}
\DeclareMathOperator{\leaksync}{\blacktriangleright_2}

%% memory update 
\DeclareMathOperator*{\updatedby}{\circledast}
\DeclareMathOperator*{\uninfupdatedby}{\circledcirc}


%% function restriction
\newcommand\restr[2]{{% we make the whole thing an ordinary symbol
  \left.\kern-\nulldelimiterspace % automatically resize the bar with \right
  #1 % the function
  \vphantom{\big|} % pretend it's a little taller at normal size
  \right|_{#2} % this is the delimiter
  }}
%% intersects
\DeclareMathOperator*{\intersects}{\dot\cap}
%% List comprehension, taken from braket
{\catcode`\|=\active
  \xdef\List{\protect\expandafter\noexpand\csname List \endcsname}
  \expandafter\gdef\csname List \endcsname#1{\left[%
     \ifx\SavedDoubleVert\relax \let\SavedDoubleVert\|\fi
     \:{\let\|\SetDoubleVert
     \mathcode`\|32768\let|\SetVert
     #1}\:\right]}
}

%% absolute/length
\providecommand{\abs}[1]{\lvert#1\rvert}
\providecommand{\len}[1]{\lvert#1\rvert}

%% from
\newcommand{\from}\leftarrow


%% Aliases for readablity
\newcommand{\Iota}I

%% checkmark and x
\newcommand{\cmark}{\ding{51}}%
\newcommand{\xmark}{\ding{55}}%

%% When the machines agree
\providecommand{\MI}{\ast}
\providecommand{\M}[1]{\MI} %can be replaced by #1 in case we decide to distinguish the machine types after all



%normal red is too bright
\definecolor{red}{HTML}{800000} 

%\usepackage{bm}


% \C is used inside figures created by inkscape, which sadly does not have good alignment support for tex.
\newcommand\C[1]{\raisebox{-0.5\height}{\texttt{#1}}} 

\title{Verifying Concurrent Operating Systems}
\author{Jonas Oberhauser}
\date{\today} % Date, can be changed to a custom date
\begin{document}
\begin{frame}
\titlepage % Print the title page as the first slide
\end{frame}


\begin{frame}
Sequential verification
\begin{center}
	\texttt{
		\begin{tabular}{l}
			i++; 
		\\ i++;
		\end{tabular}
	}
\end{center}

After execution do we have 
\[ \texttt{i} = 2 \ ? \]

Trivial
\end{frame}



\begin{frame}
In reality (e.g., on MIPS):
\begin{center}
	\texttt{
		\begin{tabular}{l}
			fetch;
			\\t1 = i;
			\\fetch;
			\\ t1 = t1 + 1; 
			\\fetch;
			\\ i = t1;
			\\
			\\ fetch;
			\\ t1 = i;
			\\ fetch;
			\\ t1 = t1 + 1; 
			\\ fetch;
			\\ i = t1;
		\end{tabular}
	}
\end{center}

After execution do we have 
\[ \texttt{i} = 2 \ ? \]

Equally Trivial
\end{frame}


\begin{frame}
Multiple threads, shared state

\begin{center}
	\texttt{
		\begin{tabular}{l||l}
			i++; & i++;
		\end{tabular}
	}
\end{center}

After execution do we have 
\[ \texttt{i} = 2 \ ? \]

\end{frame}


\begin{frame}
Interleavings:

\begin{center}
	\scalebox{.6}{
		\begin{tabular}{ccc}
			\includefig{increment_coarse12} & \hspace{4em} & \includefig{increment_coarse21} 
		\end{tabular}}
\end{center}
\end{frame}


\begin{frame}
Break up into small steps:

\begin{center}
	\texttt{
		\begin{tabular}{l||l}
			t1 = i; & t2 = i;
		\\	i = t1+1; & i = t2+1;
		\end{tabular}
	}
\end{center}

\onslide<2->
Look at all interleavings:


\begin{center}
	\scalebox{.45}{
		\begin{tabular}{ccccc}
			\includefig{increment_schedule1122} & \hspace{4em} & \includefig{increment_schedule1212} & \hspace{4em} & \includefig{increment_schedule1221}\\
			\includefig{increment_schedule2211} & \hspace{4em} & \includefig{increment_schedule2121} & \hspace{4em} & \includefig{increment_schedule2112}	
	\end{tabular}}
\end{center}

New behavior:
\[ \texttt{i} = 1 \ \lor \ \texttt{i} = 2 \]
\end{frame}




\begin{frame}
Break up into smaller steps:

\begin{center}
	\texttt{
		\begin{tabular}{l||l}
			fetch; & fetch;
			\\ t1 = i; & t2 = i;
			\\ fetch; & fetch;
			\\ t1 = t1+1; & t2 = t2+1;
			\\ fetch; & fetch;
			\\ i = t1; & i = t2;
		\end{tabular}
	}
\end{center}

New behaviors?

Look at all 924 interleavings???
\onslide<2->
\begin{center}
	\scalebox{0.6}{\includegraphics{figures/nope-im-out}}
\end{center}
\end{frame}

\begin{frame}
Step 1 to a better life: \textbf{order reduction}

\begin{enumerate} 
	\item Answers the question: 
	
	\begin{center}
		``When did we break it up enough to get all possible behaviors?''
	\end{center}
	\item Theorem: if program satisfies some software conditions, breaking it down further will not add new behaviors
\end{enumerate}
\end{frame}



\begin{frame}
Breaking up can be modeled by cooperative vs preemptive interleaving

\begin{center}
	\texttt{
		\begin{tabular}{l||l}
			t1 = i, & t2 = i,
			\\	i = t1+1; & i = t2+1;
		\end{tabular}
	}
\end{center}

\begin{description}
	\item<2->[Cooperative scheduler] will only interleave at `;' (called \textbf{interleaving points} (IP))
	
	\begin{center}
		\scalebox{.3}{
			\begin{tabular}{ccc}
				\includefig{increment_schedule1122} & \hspace{4em} & \includefig{increment_schedule2211} \\
		\end{tabular}}
	\end{center}
	\item<3->[Preemptive scheduler] will interleave at `;' and `,'
	
	\begin{center}
		\scalebox{.3}{
			\begin{tabular}{ccccc}
				\includefig{increment_schedule1122} & \hspace{4em} & \includefig{increment_schedule1212} & \hspace{4em} & \includefig{increment_schedule1221}\\
				\includefig{increment_schedule2211} & \hspace{4em} & \includefig{increment_schedule2121} & \hspace{4em} & \includefig{increment_schedule2112}	
		\end{tabular}}
	\end{center}
\end{description}

\end{frame} 



\begin{frame}
\begin{center}
	\texttt{
		\begin{tabular}{l||l}
			t1 = i, & t2 = i,
			\\	i = t1+1; & i = t2+1;
		\end{tabular}
	}
\end{center}

Observation: 
\begin{enumerate} 
	\item \texttt{t1} only accessed by Thread 1
	\item \texttt{t2} only accessed by Thread 2
	\item \texttt{i} accessed by both, concurrently
\end{enumerate}

Make this formal with ownership
\end{frame} 

\begin{frame}
\alt<2->{\alt<3->{Introduce dynamic ownership:
		\begin{center}
			\includefig{MLayoutOwnedByTwo}
	\end{center}}{Introduce dynamic ownership:
		\begin{center}
			\includefig{MLayoutOwnedAnnotated}
\end{center}}}{Basic layout: processor registers and shared memory
	\begin{center}
		\includefig{MLayoutUnowned}
\end{center}}
\end{frame} 


\begin{frame}
\textbf{Software Condition 1)}: annotate accesses based on ownership as local or shared

\begin{center}
	\includefig{MLayoutOwnedAnnotations}
\end{center}

Shared access is \textbf{linearization point} (LP)
\end{frame} 

\begin{frame}
\textbf{Software Condition 2)}: at most one LP between two IPs

Violated:
\begin{center}
	\texttt{
		\begin{tabular}{l||l}
			t1 = \{i\}, & t2 = \{i\},
			\\	\{i\} = t1+1; & \{i\} = t2+1;
		\end{tabular}
	}
\end{center}

Satisfied:
\begin{center}
\texttt{
	\begin{tabular}{l||l}
		fetch, & fetch,
		\\ t1 = \{i\}; & t2 = \{i\};
		\\ fetch, & fetch,	
		\\ t1 = t1 +1, & t2 = t2 + 1,
		\\ fetch, & fetch,
		\\	\{i\} = t1; & \{i\} = t2;
	\end{tabular}
}
\end{center}
\end{frame}


\begin{frame}
\textbf{Software Condition 3)} (technical): After reaching LP, an IP must be reachable

Violated:
\begin{center}
\texttt{
	\begin{tabular}{l}
		\{i\} = 1,
		\\ while(1) \{
		\\ \};
	\end{tabular}
}
\end{center}
\end{frame} 

\begin{frame}
Checking safety
\begin{center}
	\texttt{
		\begin{tabular}{l||l}
			fetch, & fetch,
			\\ t1 = \{i\}; & t2 = \{i\};
			\\ fetch, & fetch,	
			\\ t1 = t1 +1, & t2 = t2 + 1,
			\\ fetch, & fetch,
			\\	\{i\} = t1; & \{i\} = t2;
		\end{tabular}
	}
\end{center}
\end{frame}


\bibliography{bibliography}



\end{document}